% !PRINTED 427
\weberpage{427}
If we set, as in § 54, (4)
\[
\gamma_2(\omega) = \sqrt[3]{j(\omega)}, \quad \gamma_3(\omega) = \sqrt[2]{j(\omega)-1728},
\]
then it follows from § 54, (14):
\[
\gamma_3(i) = -\gamma_3(i), \quad \text{hence} \quad \gamma_3(i) = 0.
\]
Furthermore, by the same formula, if we set
\[
\varrho = \frac{-1+i\sqrt{3}}{2}, \qquad \varrho+1 = \frac{-1}{\varrho}
\]
then:
\[
\gamma_2(\varrho+1) = e^{-\frac{2\pi i}{3}}\gamma_2(\varrho) = \gamma_2\left(-\frac{1}{\varrho}\right) = \gamma_2(\varrho),
\]
consequently $\gamma_2(\varrho) = 0$, and we obtain
\weq{1}{\begin{aligned}
H_3(u) &= u, & j(\varrho) &= 0, \\
H_4(u) &= u-1728, & j(i) &= 1728.
\end{aligned}}

The two functions $H_3(u)$ and $H_4(u)$ therefore have integral rational numerical coefficients. This property we can now establish in general for all $H_m$ by complete induction.

First we observe that, once it has been proved that $H_m(u)$ has \gesperrt{rational coefficients}, it follows at once that the coefficients are \gesperrt{integral rational numbers} (the first $=1$). For these coefficients are integral algebraic numbers (by § 115, VI.), and consequently, if they are rational, \gesperrt{integral rational numbers}. Now we consider the largest among the values of $m$ that occur in equation (10), § 115; these are
\[
\begin{aligned}
m &= 4n, & x &= 1, & y &= 0, & k &= 1, \\
m &= 4n-1, & x &= 1, & y &= \pm 1, & k &= 2.
\end{aligned}
\]

Accordingly $F_n(u,u)$ contains the two factors $H_{4n}(u)$, $H_{4n-1}(u)^2$, and otherwise nothing but factors $H_m(u)$ in which $m<4n-1$. If we assume that it has already been proved of the latter that they have rational coefficients, it follows that
\[
H_{4n}(u)\, H_{4n-1}(u)^2 = \varPhi(u)
\]
has rational coefficients, and hence, in order to find $H_{4n-1}(u)$, one has to seek the greatest common divisor of $\varPhi(u)$ and $\varPhi'(u)$, which is done by rational calculation. But thereby $H_{4n}(u)$ too is found by rational means, and since $H_3(u)$ and $H_4(u)$ have rational coefficients, it follows in general:

VII. \gesperrt{The class functions $H_m(u)$ have integral rational numerical coefficients.}
